% Sablon pentru realizarea lucrarii de licenta, conform cu recomandarile
% din ghidul de redactare:
% - https://fmi.unibuc.ro/finalizare-studii/
% - https://drive.google.com/file/d/1xj9kZZgTkcKMJkMLRuoYRgLQ1O8CX0mv/view

% Multumiri lui Gabriel Majeri, acest sablon a fost creat pe baza
% codului sursa a lucrarii sale de licenta. 
% Codul sursa: https://github.com/GabrielMajeri/bachelors-thesis
% Website: https://www.gabrielmajeri.ro/
%
% Aceast sablon este licentiat sub Creative Commons Attribution 4.0 International License.

\documentclass[12pt, a4paper]{report}

% Suport pentru diacritice și alte simboluri
\usepackage{fontspec}

% Suport pentru mai multe limbi
\usepackage{polyglossia}

% Setează limba textului la română
\setdefaultlanguage{english}
% Am nevoie de engleză pentru rezumat
\setotherlanguages{english}

% Indentează și primul paragraf al fiecărei noi secțiuni
\SetLanguageKeys{english}{indentfirst=true}

% Suport pentru diferite stiluri de ghilimele
\usepackage{csquotes}

\DeclareQuoteStyle{english}
  {\quotedblbase}
  {\textquotedblright}
  {\guillemotleft}
  {\guillemotright}

% Utilizează biblatex pentru referințe bibliografice
\usepackage[
    maxbibnames=50,
    sorting=nty
]{biblatex}

\addbibresource{bibliography.bib}

% Setează spațiere inter-linie la 1.5
\usepackage{setspace}
\onehalfspacing

% Modificarea geometriei paginii
\usepackage{geometry}

% Include funcțiile de grafică
\usepackage{graphicx}
% Încarcă imaginile din directorul `images`
\graphicspath{{./images/}}

% Listări de cod
\usepackage{listings}

% Linkuri interactive în PDF
\usepackage[
    colorlinks,
    linkcolor={black},
    menucolor={black},
    citecolor={black},
    urlcolor={blue}
]{hyperref}

% Simboluri matematice codificate Unicode
\usepackage[warnings-off={mathtools-colon,mathtools-overbracket}]{unicode-math}

% Comenzi matematice
\usepackage{amsmath}
\usepackage{mathtools}

% Formule matematice
\newcommand{\bigO}[1]{\symcal{O}\left(#1\right)}
\DeclarePairedDelimiter\abs{\lvert}{\rvert}

% Suport pentru rezumat în două limbi
% Bazat pe https://tex.stackexchange.com/a/70818
\newenvironment{abstractpage}
  {\cleardoublepage\vspace*{\fill}\thispagestyle{empty}}
  {\vfill\cleardoublepage}
\renewenvironment{abstract}[1]
  {\bigskip\selectlanguage{#1}%
   \begin{center}\bfseries\abstractname\end{center}}
  {\par\bigskip}

% Suport pentru anexe
\usepackage{appendix}

% Stiluri diferite de headere și footere
\usepackage{fancyhdr}

% Metadate
\title{Reinforcement learning in video games}
\author{Smarandache Mihnea}

% Generează variabilele cu @
\makeatletter

\begin{document}

% Front matter
\cleardoublepage
\let\ps@plain

% Pagina de titlu
\begin{titlepage}

% Redu marginile
\newgeometry{left=2cm,right=2cm,bottom=1cm}

\begin{figure}[!htb]
    \centering
    \begin{minipage}{0.2\textwidth}
        \includegraphics[width=\linewidth]{logo-ub.png}
    \end{minipage}
    \begin{minipage}{0.5\textwidth}
        \large
        \vspace{0.2cm}
        \begin{center}
            \textbf{UNIVERSITATEA DIN BUCUREȘTI}
        \end{center}
        \vspace{0.3cm}
        \begin{center}
            \textbf{
                FACULTATEA DE \\
                MATEMATICĂ ȘI INFORMATICĂ
            }
        \end{center}
    \end{minipage}
    \begin{minipage}{0.2\textwidth}
        \includegraphics[width=\linewidth]{logo-fmi.png}
    \end{minipage}
\end{figure}

\begin{center}
\textbf{SPECIALIZAREA INFORMATICĂ}
\end{center}

\vspace{1cm}

\begin{center}
\Large \textbf{Disertație}
\end{center}

\begin{center}
\huge \textbf{\MakeUppercase{\@title}}
\end{center}

\vspace{3cm}

\begin{center}
\large \textbf{Absolvent \\ \@author}
\end{center}

\vspace{0.25cm}

\begin{center}
\large \textbf{Coordonator științific \\ Păduraru Ciprian Ionuț}
\end{center}

\vspace{2cm}

\begin{center}
\Large \textbf{București, septembrie 2023}
\end{center}
\end{titlepage}
\restoregeometry
\newgeometry{
    margin=2.5cm
}

\fancypagestyle{main}{
  \fancyhf{}
  \renewcommand\headrulewidth{0pt}
  \fancyhead[C]{}
  \fancyfoot[C]{\thepage}
}

\addtocounter{page}{1}

% Rezumatul
\begin{abstractpage}

% \begin{abstract}{english}
% Lorem ipsum dolor sit amet, consectetur adipiscing elit. Fusce vitae eros sit amet sem ornare varius. Duis eget felis eget risus posuere luctus. Integer odio metus, eleifend at nunc vitae, rutrum fermentum leo. Quisque rutrum vitae risus nec porta. Nunc eu orci euismod, ornare risus at, accumsan augue. Ut tincidunt pharetra convallis. Maecenas ut pretium ex. Morbi tellus dui, viverra quis augue at, tincidunt hendrerit orci. Lorem ipsum dolor sit amet, consectetur adipiscing elit. Aliquam quis sollicitudin nunc. Sed sollicitudin purus dapibus mi fringilla, nec tincidunt nunc eleifend. Nam ut molestie erat. Integer eros dolor, viverra quis massa at, auctor.
% \end{abstract}

% \begin{abstract}{english}
% Lorem ipsum dolor sit amet, consectetur adipiscing elit. Fusce vitae eros sit amet sem ornare varius. Duis eget felis eget risus posuere luctus. Integer odio metus, eleifend at nunc vitae, rutrum fermentum leo. Quisque rutrum vitae risus nec porta. Nunc eu orci euismod, ornare risus at, accumsan augue. Ut tincidunt pharetra convallis. Maecenas ut pretium ex. Morbi tellus dui, viverra quis augue at, tincidunt hendrerit orci. Lorem ipsum dolor sit amet, consectetur adipiscing elit. Aliquam quis sollicitudin nunc. Sed sollicitudin purus dapibus mi fringilla, nec tincidunt nunc eleifend. Nam ut molestie erat. Integer eros dolor, viverra quis massa at, auctor.
% \end{abstract}

\begin{abstract}{english}
    Video games have always been a domain where technological innovations were always important and necessay for its advancement. One of the more recent advancements in technology was the brining of the fields of machine learning and deep learning into the mainstream. These two could prove to be real game changers in the industry of video game development, assisting in the development of AI systems. The role of this paper is to implement simple AI behaviours commonly used in video games using Reinforcement Learning, and compare the results with a more classical approach.
\end{abstract}

% jocurile video au fost 

\end{abstractpage}

\tableofcontents

% Main matter
\cleardoublepage
\pagestyle{main}
\let\ps@plain\ps@main

\chapter{Introduction}

The video games industry is one that has been around since 1970 and over the years it became one of the most successful forms of entertainment. As the years went by, the complexity of the games also increased, requiring more time to develop and more manpower. In the early days, mainstream titles would have smaller teams, usually comprised of 10 to 20 people, however, as time went by, so did the teams increase in size, with studios having over hundreds of people working on a single game.

One important aspect of video games is the AI. It is a tough challenge for a developer to create and implement an AI to perform actions in such a way that feels \emph{natural} and not scripted, so that the player's immersion is not broken. As the video games evolved, so did the AI, becoming more complex and having more ways to interact with the player and the environment.

A solution for this problem would be to implement Reinforcement Learning techniques to train the AI to perform the required tasks. Some advantages of this approach would be that the developer would need to create a set of rules for the AI, and then reward or punish the it for its taken actions. Another advantage would be that it would be simple to implement this behaviour, if the actions that the AI perform are similar to the actions that the player performs since they have already been implemented. However, this approach also has some negative aspects. For starters, this training could take a long period of time, depending on the complexity of this problem and the allotted resources for the training process. Another disadvantage would be that the even if the developers do not have to manually implement the AI logic, they still have to create punishments and rewards for the training agent, and also have to balance these rewards to reach an acceptable result, which could take time due to trial-and-error.

% industria e una foarte profitabila, e o chestie mainstream
% un aspect important al jocurilor este AI-ul, care trebuie sa fie destul de bun incat sa ofere o provocare jucatorilor, dar sa nu fie unfair
% datorita faptului ca jocurile devin din ce in ce mai complexe, e din ce in ce mai greu sa dezvolti un AI 
% dezvoltarea cu RL are avantajul ca doar ii dai anumite reguli unui obiect si el invata in functie de ele, poate fi oarecum mai impredictable
% dezavantajul e ca trebuie sa vezi exact cum il pedepsesti sau recompensezi, plus ca dureaza sa il antrenezi
%

\chapter{Preliminary}


\section{Reinforcement Learning \& PPO}

% TODO: aici vreo 3-4 pagini max nu vreau sa scriu prea mult dar tre sa ma bazez pe referinte
According to \cite{Fran_ois_Lavet_2018_rlintro}, Reinforcement Learning can be defined as \enquote{the area of machine learning that deals with sequential decision-making}; it is a method of training an agent to learn a specific behaviour by interacting and observing the environment that it exists inside of. With each action performed by the agent which affects the environment, it can be rewarded or punished so that with each iteration it can adapt its behaviour to be slightly more similar to the desired behaviour. 

\cite{Fran_ois_Lavet_2018_rlintro} defines the general Reinforcement Learning problem as a \enquote{discrete time stochastic control process where an agent interacts with its environment in the following way: the agent starts, in a given state within its environment $s_0 \in \gS$, by gathering an initial observation $\omega_0 \in \Omega$. At each time step $t$, the agent has to take an action $a_t \in \gA$. [\dots]. it follows three consequences: (i) the agent obtains a reward $r_t \in \gR$, (ii) the state transitions to $s_{t+1} \in \gS$, and (iii) the agent obtains an observation $\omega_{t+1} \in \Omega$}. Because of this, Reinforcement Learning can be modeled as a Markov Decision Process (MDP). According to \cite{Bellman_MDP}, a MDP is a 4-tuple $(\gS, \gA, T, R)$ where:
\begin{itemize}
    \item [$\gS$]: is the state space
    \item [$\gA$]: is the action space
    \item [$T$]: is the transition function; it is defined as $T: \gS \times \gA \times \gS \rightarrow [0, 1]$, and represents the probability that an action $a \in \gA$ will lead from state $s \in \gS$ to state $s' \in gS$
    \item [$R$]: is the reward function; it is defined as $R: \gS \times \gA \times \gS \rightarrow \gR$, and represents the reward obtained when moving from state state $s \in \gS$ to state $s' \in gS$ because of action $a \in \gA$
\end{itemize}

Another notion of Reinforcement Learning is the policy, which shows how an agent selects the action that it will take. It can be defined as $\pi: \gA \times \gS \rightarrow [0, 1]$, where $\pi(a, s)$ is the probability of action $a$ being chosen while in state $s$.

To solve a Reinforcement Learning problem there are multiple ways, such as model-based methods, value-based methods, policy gradient methods, etc. The ones that are important for this paper are the methods with policy gradients, since those are the ones that will be used. Policy gradient methods learn a parameterized policy and optimize it directly. The policy can be defined as 
\begin{equation} \label{policy_eq:1}
    \pi_{\theta} (a | s) = Pr\{a_t = a | s_t = s, \theta_{t} = \theta\}
\end{equation}
where $\theta \in \sR^{d'}$ is the policy's parameter vector.




According to \cite{schulman2017ppo} \enquote{Policy gradient methods work by computing an estimator of the policy gradient and plugging it into a stochastic gradient ascent algorithm. The most commonly used gradient estimator has the form 

\begin{equation}
    \hat{g} = \hat{\mathbb{E}}_{t} [\nabla_{\theta} log \pi_{\theta}(a_t | s_t) \hat{A}_{t}]    
\end{equation}
where $\pi_{\theta}$ is a stochastic policy and $\hat{A}_{t}$ is an estimator of the advantage function at timestep $t$.}

$\hat{\mathbb{E}}_{t}$ is the expectation which is used in an algorithm that alternates between optimization and sampling, and which indicates the empirical average over a finite batch of samples. Implementations utilizing automatic differentiation software operate by creating an objective function, the gradient of which serves as the policy gradient estimator. The estimator $\hat{g}$ is derived through the differentiation of the objective

\begin{equation}
    L^{PG} (\theta) = \hat{\mathbb{E}}_{t} [log \pi_{\theta}(a_t | s_t) \hat{A}_{t}]
\end{equation}

\paragraph{}
Proximal policy optimization (PPO) works by sampling the data obtained from interacting with the environment and then optimizes the objective function with stochastic gradient ascent

\begin{equation}
    r_t(\theta) = \frac{\pi_{\theta} (a_t | s_t)}{\pi_{\theta_{\text{old}}} (a_t | s_t)}
\end{equation}

\begin{equation}
    L(\theta) = \hat{\mathbb{E}}_{t} [min(r_t(\theta) \hat{A}_{t}, clip(r_t(\theta), 1-\epsilon, 1+\epsilon) \hat{A}_{t})]
\end{equation}
where $r_t(\theta)$ is the probability ratio. It can be seen that the probability function is clipped in order for the surrogate objective to be modified.

The PPO algorithm that uses fixed-length trajectory segments as defined in \cite{schulman2017ppo} can be seen in Alogrithm \ref{alg:ppo}

\begin{algorithm}
    \caption{PPO, Actor-Critic Style}\label{alg:ppo}
    \begin{algorithmic}
    \For{iteration $ = 1, 2, \dots $}
        \For{actor $ = 1, 2, \dots, N$}
            \State Run policy $\pi_{\theta_{\text{old}}}$ in environment for $T$ timesteps
            \State Compute advantage estimates $\hat{A}_{1}, \dots, \hat{A}_{T}$
        \EndFor
        \State Optimize surrogate $L$ wrt $\theta$, with $K$ epochs and minibatch size $M \leq NT$
        \State $\theta_{\text{old}} \gets \theta$
    \EndFor
    \end{algorithmic}
\end{algorithm}

At each iteration, each of the $N$ actors collect $T$ timesteps of data, then the surrogate loss is constructed on the $NT$ timesteps of data and optimized witha  minibatch SGD for $K$ epochs.

% sa mai copiez niste cacat pe aici

\section{Unity \& ML-Agents Toolkit}

The environment where the solution will be implemented will be a Unity project consisting of a small level that has different structures an objects, and which contains tanks that can move, shoot, and fight eachother. The project also contains an AI implementation that is based on behaviour trees and genetic algorithms and is described in more detail in \cite{paduraru2019automatic}. The project can be found at \url{https://github.com/AGAPIA/BTreeGeneticFramework}.

The ML-Agents Toolkit is defined as \enquote{an open source project which enables researchers and developers to create simulated environments using the Unity Editor and interact with them via a Python API. The toolkit provides the ML-Agents SDK which contains all functionality necessary to define environments within the Unity Editor along with the core C\# scripts to build a learning pipeline} \cite{juliani2020unityai}. The toolkit also provides implementation in PyTorch for most state-of-the-art reinforcement learning algorithms, including PPO, and several example environments to quickly test the toolkit.

There are 3 entities provided by the ML-Agents SDK:
\begin{itemize}
    \item sensors, which are used to collect the observations from the environment
    \item agents, which are components to add the Agent functionality to a Unity GameObject, which enables them to collect observations, decide which action to take by using the observations and execute them, and receive rewards
    \item academy, which can be defined as a singleton instance that manages agent training and decision making \cite{mlagents_api_docs}
\end{itemize}

The observations that are collected, the actions that should be taken by the agent, and the rewards received by the agent can be implemented manually using the provided \emph{Agent} class.

The ML-Agents Toolkit also has a Python package, \emph{mlagents}, which according to \cite{mlagents_pypi}, provides a set of reinforcement and imitation learning algorithms designed to be used with Unity environments. This package will be used to effectively run the training process, by communicating with the Unity application via a gRPC communication protocol.

% sa bag poze, sa zic ca mlagents imi ofera moduri de a scripta ce obs fac ce pedepse si recompense dau pt agent, modalitate de a antrena si a face inferenta cu ce am antrenat probabil sa iau ceva si de la \cite{juliani2020unityai}
\chapter{Training TODO: change name, kinda cringe}
\section{Reaching a static target}
\subsection{The Problem}
A common task for an AI in a game is to reach a a given destination. According to \cite{alonso2020deeplearningnavigation} the common way to achieve this goal is by using \emph{navigation meshes} (NavMeshes). These are represented as graphs, and their nodes represent surfaces that can be traversed. Afterwards, algorithms such as \emph{A*} can be used to find the fastest way for an agent to go from point A to point B. However, these \emph{NavMeshes} require to be \emph{baked} (TODO explicatie baked) in advance, so updating them in real time can prove to be a challenge depending on several factors such as:
\begin{itemize}
    \item the game engine used
    \item the complexity of the game environment
    \item the computational cost associated with recreating in real time these meshes
\end{itemize}

A proposed solution for this problem would be to use AI agents that have been trained using deep learning methods, in such a manner that they would be independent from changes to the environment.

\subsection{Implementing the solution}
\paragraph{}
First, the action space for the agent is defined: due to the fact that we only need to move the agent to a given position, the action space consists of moving the agent forward or backward, and rotating it. Because the agent is supposed to be controlled by a controller, its movement input is defined as a real number in the $[-1, 1]$ interval for both $X$ and $Y$ axes. However, a discrete action space will be used instead of a continuous one to reduce the difficulty of learning, and also because there is no need for the agent to have movements that are so precise. For each axis the action space will be represented by the discrete space: $\{-1, -0.5, 0, 0.5, 1\}$.

\paragraph{}
To make decisions, the agent will need to make several observations about its surroundings. Firstly, it will need to know how close it is to certain objects in the environment. To accomplish this, several \emph{raycasts} will be used to measure the distance from the agent, similar to a \emph{LIDAR} (TODO: add ceva despre lidar si un paper). The raycasts start from the center of the agent and are spread in such a way that the angle between 2 consecutive rays is equal for any 2 consecutive rays (TODO: imagine cu rays, si probabil sa gasesc un termen mai ok). The observation will contain the distance until the rays hit an object. Also, the index of the layer of the hit object will be included in the observations, so that the agent will differentiate between regular environment objects and more important objects, such as other players, enemies, etc. For this implementation, 32 rays were used.

Other observations that are made are the agent's position in space and also that of the target. This observation is included so that the agent can learn how movement brings it closer or further to the target. The next observation is the agent's forward vector so it can know in which direction it is moving. The optimal direction that the agent should take is also observed and obtained by computing the vector difference between the target's position and the agent's position. To know how much to adjust its trajectory, the angle between the agent's forward vector and the target is observed; the angle is signed so that the agent can learn to adjust its trajectory to the left or to the right. The distance to the target is added to the observations so that the agent can learn that when the distance is getting smaller it is rewarded. The agent's normalized velocity vector is observed to show in which direction it is moving based on the given input. The angle between the agent's velocity vector and its forward vector is observed to tell if the agent is moving forwards or backwards. Finally, the agent's velocity magnitude is observed so the agent can know if it is moving or standing still. 

In summary the following observations are being made:
\begin{itemize}
    \item distance for each raycast until it hits an object
    \item layer index of object hit by the raycast
    \item spatial postion of the agent
    \item spatial position of the target
    \item agent's forward vector
    \item optimal direction of the agent
    \item signed angle between the agent's forward vector and the target
    \item distance from the target
    \item agent's normalized velocity vector
    \item angle between the agent's velocity vector and its forward vector
    \item agent's velocity magnitude
\end{itemize}

\paragraph{}
In order for the agent to be able to learn to solve a specific problem, in this case, to reach a target, it must be rewarded or punished according to the actions that were taken. It is important to be mindful of the rewards that are given to the agent because only giving rewards and not punishing it can lead the agent to learn a behaviour that will maximize its reward, but will not be able to solve the problem, as it will be described shortly.

To begin, the first reward that was implemented, was the reward for reaching the objective, which is the goal of the agent and should also be a substantial reward. The other rewards that ware added were based on the direction of the movement and how close to the target it was, the reward increasing in value if the agent was closer to the objective (\ref{reward:1}) and a reward if the raycasts are hitting the destination object, which also increases in value if the agent is closer to the target (\ref{reward:2}). 

\begin{equation} \label{reward:1}
    R = \frac{(1 - \frac{\alpha}{180}) \cdot r}{d}
\end{equation}
where:
\begin{itemize}
    \item [$R$]: is the total reward
    \item [$\alpha$]: is the angle between the agent's current direction and the optimal direction
    \item [$r$]: is the reward that is obtained if the agent has the optimal direction
    \item [$d$]: is the distance from the agent to the objective
\end{itemize}

\begin{equation} \label{reward:2}
    R = \frac{r}{d_{i}}
\end{equation}
where:
\begin{itemize}
    \item [$R$]: is the total reward
    \item [$r$]: is the reward obtained if the distance between the agent and the objective is minimum
    \item [$d_{i}$]: is the obtained distance by the raycast $i$ from the agent to the object
\end{itemize}

However, these three rewards are not enough for the agent: through learning experiments it was observed that the agent was performing poorly: it would get stuck trying to move through a wall, it would never reach the objective and just spin around, or in some cases, it would just stand still.

To solve these problems, several punishments were implemented to correct the behavior of the agent. To prevent the agent for not moving, a constant penalty was added, to incentivise the agent to move towards the reward, so that it will receive a reward. Also, to comabat standing still, if the agent does not move, it receives an additional penalty, and if it does not move for more than 100 time steps, it receives a huge penalty and the episode ends. 

To stop the agent from trying to pass thorugh walls, a penalty is added if the raycasts that hit walls or other objects in the environment, have the distance to the hit object be a smaller than a given number. This penalty also increases the closer the agent gets to a wall (\ref{punishment:1}). 

\begin{equation} \label{punishment:1}
    R = \frac{p}{d_{i}}, \text{if } d_{i} < d
\end{equation}
where:
\begin{itemize}
    \item [$R$]: is the total reward
    \item [$p$]: is the penalty obtained if the distance between the agent and the wall/environmental object is minimum
    \item [$d_{i}$]: is the obtained distance by the raycast $i$ from the agent to the object
    \item [$d$]: is the maximum distance for which the penalty is applied
\end{itemize}

Another penalty was added if the agent is moving away from objective, and as before, it increases the further away it gets from the objective (\ref{punishment:2}). 

\begin{equation} \label{punishment:2}
    R = \frac{\alpha}{180} \cdot p
\end{equation}
where:
\begin{itemize}
    \item [$R$]: is the total reward
    \item [$p$]: is the penalty obtained if the distance between the agent and the wall/environmental object is minimum
    \item [$\alpha$]: is the angle between the agent's current direction and the optimal direction
\end{itemize}

Through training, two undesirable behaviours were observed: it was observed that the agent would sometimes make sudden jerky movments, trying to change its direction and immediatly returning to its previous trajectory, and that the agent would learn to drive backwards. To fix the first problem, a penalty was added if the agent would change its direction (\ref{punishment:3}), and to fix the second one, a penalty was added if the tank was moving backwards (\ref{punishment:4}).

\begin{equation} \label{punishment:3}
    R = \frac{\alpha}{180} \cdot p
\end{equation}
where:
\begin{itemize}
    \item [$R$]: is the total reward
    \item [$p$]: is the penalty obtained for changing the movement direction
    \item [$\alpha$]: is the angle between the agent's current direction and its previous direction
\end{itemize}

\begin{equation} \label{punishment:4}
    R = \frac{\alpha}{180} \cdot p, \text{if } \alpha > 90
\end{equation}
where:
\begin{itemize}
    \item [$R$]: is the total reward
    \item [$p$]: is the penalty obtained for changing the movement direction
    \item [$\alpha$]: is the angle between the agent's current direction and its forward vector
\end{itemize}

In summary, the used rewards and penalites and their values can be seen in table \ref{reward_punish_table:1}

\begin{table}
    \begin{tabular}{|| m{15em} | m{4em} | m{15em} ||}
    \hline \hline
    Name & Value & Notes \\ \hline \hline
    Reach Objective Reward & 10 &  \\ \hline
    Move Towards Objective Reward & 0.001 & is scaled by the distance between agent and objective \\ \hline
    Raycast Touches Objective Reward & 0.001 & is scaled by the distance between agent and objective \\ \hline
    Constant Penalty & -0.005 & is applied at each time step \\ \hline
    Not Moving Penalty & -0.05 &  \\ \hline
    Not Moving For 100 Steps Penalty & -50 &  \\ \hline
    Moving Towards Wall Penalty & -0.002 & is scaled by the distance between agent and object \\ \hline
    Moving Away From Objective Penalty & -0.025 & is scaled by the distance between agent and objective \\ \hline
    Sudden Movement Penalty & -0.01 & is scaled by the angle between the agent's current direction and its previous one \\ \hline
    Moving Backwards Penalty & -0.005 &  \\ \hline \hline
    \end{tabular}
    \caption{Rewards and Penalites}
    \label{reward_punish_table:1}
\end{table}

\subsection{Training}

Each training session consisted of $10^7$ steps, and 30 agents were trained in parallel. The training times were between 3-4 hours. The agents are supposed to learn to reach an objective which appears in one of 17 predefined positions on the map. Once the agent reaches the objective, it is moved in another location chosen randomly. Each separate training instance uses the same seed for the random number generator, so that the objectives will appear in the same order for each of the training instances.

In the initial training sessions, the agents were unable to learn to reach the objective, becoming stuck rotating in a circle. A proposed solution was to keep the objective in the same position until the agent reaches it 30 times. After reaching the objective 30 times, the objective would start appearing in the random predefined locations. This was done to possibly kickstart the agent's learning process, but the approach failed, the agent being unable to learn to reach the objective. Another approach was to increase the neural network's size, however this approach prove unsuccessful as well. The approach that worked was to remove the agent's postion and the objective's position from the observations.

TODO:

Despre ce sa scriu:
- ce observatii am ales
- ce recompense si pedepse ii dau
- faptu ca initial nu a mers cand am bagat pozitiile
- am 30 de agenti care se antreneaza in paralel
- la training sa zic despre configuratiile folosite si sa fac grafice cu mean rewards
- ca i-am bagat
\section{Reaching a moving target}

TODO:
- sa zic ca am incercat sa ii bag memorie recurenta si nu a mers
\section{Shooting a moving target}

\nocite{*} % printeaza si referintele care nu au fost folosite in text
\printbibliography[heading=bibintoc]

\end{document}